% !TeX root = ../../elegantbook-cn.tex

\section{入题思路}
根据以往的教学经验，不少同学在面对解析几何解答题时容易陷入思路混乱的困境：将直线方程与椭圆方程联立，并利用韦达定理得到两根关系后，便不知道下一步该做什么；还有一些同学虽然能够找到初步思路，但当题目中的动点和直线逐渐增多时，往往会引入大量未知数，却无法理清它们之间的关系，最终无从下手。因此，建立清晰、连贯的解题思路，是解决解析几何解答题的关键。

为了快速理清题目中各个点、各条直线以及相关图形之间的位置关系，在解答每一道题之前，都应先根据题意画出草图。草图不必十分精确，但应准确反映题目中图形的先后顺序，为后续选择未知数、设点设线和确定运算方向提供帮助。
\subsection{动点类问题}
顾名思义，动点类问题即由一个动点（动点可在椭圆、直线等曲线上）位置变化引发的求值、范围、证明类问题。


\subsubsection*{入题思路}
设动点坐标$P(x_0, y_0)$，同时写出$x_0, y_0$的约束关系，把题目中所有与$P(x_0, y_0)$相关点坐标、弦长、面积、向量、夹角等条件用$x_0, y_0$表示，最后用$x_0, y_0$的约束关系进行化简消元、完成求值。

\begin{remark}
    此类问题不一定需要联立直线与曲线使用韦达定理 \\
    如果题目中有跟$P(x_0, y_0)$完全无关的变量，设新的变量继续翻译条件即可
\end{remark}

\subsubsection*{变量化简}
运算全程要主动利用$x_0, y_0$的约束条件，随时化简、减少变量。
如果是直线关系，可以直接用$x_0$表示$y_0$，则可以全过程只带有$x_0$一个未知数;
如果是二次曲线关系（如椭圆），当出现$x_0^2, y_0^2$时，优先利用曲线方程代换降次，减少未知数。

\begin{remark}
思考：运算结果同时出现$x_0^2, y_0^2$，那应该代换哪一个？
\end{remark}

\begin{example}
（2026海淀一模）已知椭圆$E:\frac{x^2}{a^2}+\frac{y^2}{b^2}=1(a > b > 0)$的离心率为$\frac{\sqrt{3}}{2}$，其左、右焦点和上顶点构成的三角形面积为$3\sqrt{3}$ \\
（1）求椭圆$E$的方程；\\
（2）设$O$为坐标原点，$M(x_0, y_0)$为第一象限内$E$上的动点，点$N$在直线$y=-2$上，且$OM\bot ON$，过点$A(0,3)$作$MN$的垂线交$OM$于$B$ ，求$\vert AB \vert$的值。
\end{example}
\begin{property}
\vspace*{7\baselineskip}
\end{property}
\newpage
\begin{example}
（2025海淀一模）
已知椭圆$W:\frac{x^2}{a^2}+\frac{y^2}{b^2}=1(a > b > 0)$，$A, B$分别是$W$的左右顶点，$C$是$W$的上顶点，$\Delta ABC$的面积为2，且$\vert AC \vert = \sqrt{5}$ \\
（1）求椭圆$W$的方程及长轴长；\\
（2）已知点$M(2,1)$，点$P$在直线$AC$上，设直线$PM$与$x$轴交于点$E$，直线$BP$与直线$EC$交于点$Q$，判断点$Q$是否在椭圆$W$上，并说明理由。
\end{example}
\begin{property}
\vspace*{15\baselineskip}
\end{property}
\begin{example}
（2024海淀一模）已知椭圆$G: x^2 + my^2 =m$的离心率为$\frac{\sqrt{2}}{2}$，$A_1, A_2$分别是$G$的左、右顶点，$F$是$G$的右焦点\\
（1）求$m$的值及点$F$的坐标；\\
（2）设$P$是椭圆$G$上异于顶点的动点，点$Q$在直线$x = 2$上，且$PF \bot FQ$，直线$PQ$与$x$轴交于点$M$，比较$\vert MP \vert ^2$与$\vert MA_1\vert \cdot \vert MA_2 \vert$的大小

\end{example}
\begin{property}
\vspace*{15\baselineskip}
\end{property}
\newpage
\begin{example}
（2023北京高考）已知椭圆 $E:\frac{x^2}{a^2}+\frac{y^2}{b^2}=1\ (a>b>0)$ 的离心率为 $\frac{\sqrt{5}}{3}$，$A,C$ 分别是 $E$ 的上、下顶点，$B,D$ 分别是 $E$ 的左、右顶点，$\lvert AC\rvert=4$.\\
(1) 求 $E$ 的方程；\\
(2) 设 $P$ 为第一象限内 $E$ 上的动点，直线 $PD$ 与直线 $BC$ 交于点 $M$，直线 $PA$ 与直线 $y=-2$ 交于点 $N$. 求证：$MN\parallel CD$.
\end{example}
\begin{property}
\vspace*{15\baselineskip}
\end{property}
\subsection{直线类问题}
顾名思义，直线类问题是一条直线与椭圆相交，由这条直线引出的求值、范围、证明类问题。

\subsubsection*{入题思路}
依据题干条件合理设定直线方程，设直线与椭圆的两个交点坐标分别为\((x_1, y_1)、(x_2, y_2)\)；联立直线与椭圆方程，整理得出关于$x$的一元二次方程，借助韦达定理写出两根关系：\(x_1+x_2=-\dfrac{b}{a}\)，\(x_1x_2=\dfrac{c}{a}\)。
随后将题干给出的点坐标、弦长、面积、向量、夹角等各类几何条件，全部依托\(x_1,x_2,y_1,y_2\)列式表达；再根据直线解析式，把$y_1,y_2$统一替换为含$x_1,x_2$的式子，此时式子中仅剩余$x_1,x_2$两个变量，最终代入韦达定理化简消元，完成计算证明求解。

当反设直线即$x = ny + t$时，联立直线与椭圆方程，整理得出关于$y$的一元二次方程，借助韦达定理写出两根关系：\(y_1+y_2=-\dfrac{b}{a}\)，\(y_1y_2=\dfrac{c}{a}\)，再根据直线解析式，把$x_1,x_2$统一替换为含$y_1,y_2$的式子，此时式子中仅剩余$y_1,y_2$两个变量。

\begin{remark}
    题目中若出现了与直线无关的条件，设新的变量即可
\end{remark}
\subsubsection*{常见几种设线的方法}
\noindent（1）点斜式（最常用）：适用于直线过定点，如过点$(1, 3)$的直线可以设为$y - 3 = k(x - 1)$，其中$k$为直线斜率；需要额外讨论直线斜率不存在的情况，即直线为$x = 1$的情况。\\
（2）斜截式：适用于一般直线，如$y = kx + m$，其中$k$为直线斜率；需要额外讨论直线斜率不存在的情况。\\
（3）横截式：适用于直线过$x$轴定点，如过点$(0,1)$的直线可以设为$x = ny + 1$，其中$n$为直线斜率的倒数；需要额外讨论直线斜率为0的情况，即直线为$y = 0$的情况。
\begin{example}
（2025东城二模）已知椭圆 $E:\frac{x^2}{a^2}+\frac{y^2}{b^2}=1\ (a>b>0)$ 的一个顶点为 $A(0,1)$，且过点 $\left(1,\frac{\sqrt{3}}{2}\right)$\\
(1) 求椭圆 $E$ 的方程及焦距；\\
(2) 过点 $\left(0,2\right)$ 的直线与椭圆 $E$ 交于不同的两点 $B,C$。直线 $AB,AC$ 的斜率分别记为 $k_1$ 与 $k_2$，当 $k_1+k_2=\frac{2}{3}$ 时，求 $\triangle ABC$ 的面积
\end{example}
\begin{property}
\vspace*{17\baselineskip}
\end{property}
\begin{example}
（2026丰台一模）已知椭圆 $E:\frac{x^2}{a^2}+\frac{y^2}{b^2}=1\ (a>b>0)$，椭圆 $E$ 上一点 $A(2,1)$ 到两焦点的距离之和为 $2\sqrt{6}$\\
(1) 求椭圆 $E$ 的方程和离心率；\\
(2) 过点 $P(2,3)$ 的直线 $l$ 与椭圆 $E$ 交于不同的两点 $B,C$。若直线 $AB$ 与直线 $AC$ 的斜率之和为 $0$，求直线 $l$ 的方程
\end{example}
\begin{property}
\end{property}
% \vspace*{17\baselineskip}
\newpage
\begin{example} % 直线，直线相交求坐标，求最值
（2024海淀期末）已知椭圆 $E:\frac{x^2}{a^2}+\frac{y^2}{b^2}=1\ (a>b>0)$ 过点 $A(3,0)$，焦距为 $2\sqrt{5}$.\\
（1）求椭圆 $E$ 的方程，并求其短轴长；\\
（2）过点 $P(1,0)$ 且不与 $x$ 轴重合的直线 $l$ 交椭圆 $E$ 于两点 $C,D$，连接 $CO$ 并延长交椭圆 $E$ 于点 $M$，直线 $AM$ 与 $l$ 交于点 $N$，$Q$ 为 $OD$ 的中点，其中 $O$ 为原点. 设直线 $NQ$ 的斜率为 $k$，求 $k$ 的最大值.
\end{example}
\begin{property}
\vspace*{15\baselineskip}
\end{property}

\begin{example} % 直线，基础翻译
（2024朝阳期末）已知椭圆 $E:\frac{x^2}{a^2}+\frac{y^2}{b^2}=1\ (a>b>0)$ 的左顶点为 $A$，上顶点为 $B$，原点 $O$ 到直线 $AB$ 的距离为 $\frac{2\sqrt{5}}{5}$，$\triangle AOB$ 的面积为 $1$.\\
（1）求椭圆 $E$ 的方程；\\
（2）过点 $P(-2,1)$ 的直线 $l$ 与椭圆 $E$ 交于不同的两点 $C,D$，过点 $C$ 作 $x$ 轴的垂线，分别与直线 $AD,AB$ 交于点 $M,N$. 判断点 $N$ 是否为线段 $CM$ 的中点，说明理由.
\end{example}
\begin{property}
\vspace*{15\baselineskip}
\end{property}



\subsection{设动点和设直线的本质区别\&不要被题目的描述给误导}

设动点的本质原理：以动点坐标为核心变量，将几何关系统一表示为关于动点坐标的代数关系，再利用动点自身的约束完成求解或证明。在动点问题中，部分题目的主动点并没有明确给出，这时应该理清题目的逻辑确定哪个是主动点，再设动点坐标完成求解。

设直线的本质原理：当椭圆上的两个动点之间存在整体约束，例如两点所在的直线经过某个定点，或者该直线满足某种性质时，若只设其中一个点，通常难以简洁地表示另一个点的坐标。此时可以设两点所在的直线为$y=kx+m$，引入斜率 $k$、截距 $m$等参数，将参数注入到两个点坐标中，让参数代替两个点参与运算，而在参数和点坐标中起到桥梁作用的就是韦达定理。
理论上，联立直线与椭圆后，可以利用一元二次方程的求根公式求出两个交点的坐标（用 $k,m$）表示，但所得结果通常较为复杂，且含有难以化简的根式。因此采用“设而不求”的思想：设直线与椭圆的交点为$(x_1, y_1), (x_2, y_2)$，不分别求出两个交点，而是将待求量或待证关系转化为关于$x_1+x_2,x_1x_2$的式子，再利用韦达定理进行整体代换，从而得到关于参数$k, m$的等式关系，从而避开复杂的求根过程，降低运算量。

思考：根据下面的描述选择设动点还是设直线
椭圆$\frac{x^2}{12} + \frac{y^2}{9} = 1$上有两点$A(0, 3), B(3, \frac{3}{2})$，过点$B$的直线$l$交椭圆于点$P$，若$S_{\triangle ABP} = 9$，求$l$的直线方程

已知椭圆$\frac{x^2}{4} + \frac{y^2}{3} = 1$，$A$为椭圆的左顶点，$B(0, 1)$，椭圆上有一动点$M$，连接$BM$交椭圆于点$C$，求$\triangle ABC$面积的最大值


如果椭圆上有两个点，设了其中一个点坐标之后，另一个点坐标可以快速、简洁表示成第一个点坐标的相关形式，这时候设点还是设线？
如果椭圆上


\begin{example} % 直线或者动点引入，三角形面积的转换（千锤法）
（2026海淀二模）椭圆 $E:\frac{x^2}{a^2}+\frac{y^2}{b^2}=1\ (a>b>0)$ 的下顶点为 $A(0,-1)$，左焦点到点 $A$ 的距离为 $\sqrt{3}$。\\
（1）求椭圆 $E$ 的方程及离心率；\\
（2）经过点 $A$ 的直线与椭圆的另一个交点为 $D$（位于 $x$ 轴上方），与 $x$ 轴的交点为 $M$。点 $D$ 关于 $x$ 轴的对称点为 $E$。过点 $D$ 作与 $x$ 轴平行的直线交直线 $AE$ 于点 $N$。设$\triangle DEN$与$\triangle AEM$的面积分别为 $S_{\triangle DEN}$ 和 $S_{\triangle AEM}$，若 $S_{\triangle DEN}=3S_{\triangle AEM}$，求直线 $AD$ 的斜率。
\end{example}

\begin{example} % 动点或者直线引入，三角形面积
（2025丰台二模）已知椭圆 $E:\frac{x^2}{a^2}+\frac{y^2}{b^2}=1\ (a>b>0)$ 的左顶点为 $A(-2,0)$，焦距为 $2\sqrt{2}$。\\
（1）求椭圆 $E$ 的方程；\\
（2）设 $O$ 为原点，过点 $A$ 且斜率为 $k$ 的直线 $l$ 与椭圆 $E$ 的另一个交点为 $T$，线段 $AT$ 的垂直平分线与 $x$ 轴交于点 $M$，与 $y$ 轴交于点 $N$。过点 $P(1,0)$ 且与 $l$ 平行的直线与 $y$ 轴交于点 $Q$。若 $\triangle OMN$ 与 $\triangle NPQ$ 的面积之比为 $4:3$，求 $k$ 的值。
\end{example}

% 多个动点，依照逻辑的先后设动点

\begin{example} % 入题思路：点M，点N，直线AM，直线AN
（2025丰台期末）已知椭圆 $C:\frac{x^{2}}{a^{2}}+\frac{y^{2}}{b^{2}}=1\ (a>b>0)$ 的上顶点为 $A(0,1)$，离心率为 $\frac{\sqrt{3}}{2}$.\\
(1) 求椭圆 $C$ 的方程；\\
(2) 设 $B$ 为椭圆 $C$ 的下顶点，动点 $M$ 到坐标原点 $O$ 的距离等于 $1$ ($M$ 与 $A,B$ 不重合)，直线 $AM$ 与椭圆 $C$ 的另一个交点为 $N$. 记直线 $BM, BN$ 的斜率分别为 $k_{1}, k_{2}$，问：是否存在常数 $\lambda$，使得 $k_{1}+\lambda k_{2}=0$ 恒成立？若存在，求出 $\lambda$ 的值；若不存在，说明理由.
\end{example}



\begin{example} % 动点引入（入题思路比较难），垂直的翻译，弦长公式
（2021海淀二模）椭圆 $C:\frac{x^2}{a^2}+\frac{y^2}{b^2}=1\ (a>b>0)$ 的左、右焦点分别为 $F_1,F_2$，$E$ 是椭圆 $C$ 上一点，且 $\lvert F_1F_2\rvert=2$，$\lvert EF_1\rvert+\lvert EF_2\rvert=4$。\\
（1）求椭圆 $C$ 的方程；\\
（2）$M,N$ 是 $y$ 轴上的两个动点（点 $M$ 与点 $E$ 位于 $x$ 轴的两侧），$\angle MF_1N=\angle MEN=90^\circ$，直线 $EM$ 交 $x$ 轴于点 $P$，求 $\frac{\lvert EP\rvert}{\lvert PM\rvert}$ 的值。
\end{example}


% 题目有误导，或者是入题思路比较困难

\begin{example} % 动点（题目误导），韦达定理求坐标，定值定点（特殊情况先得答案）
（2025朝阳期末）已知椭圆 $E:\frac{x^2}{a^2}+\frac{y^2}{b^2}=1\ (a>b>0)$ 的离心率为 $\frac{\sqrt{6}}{3}$，右顶点为 $(2,0)$.\\
（1）求椭圆 $E$ 的方程；\\
（2）过原点 $O$ 且与 $y$ 轴不重合的直线 $l$ 与椭圆 $E$ 交于 $M,N$ 两点. 已知点 $P(0,2)$，直线 $PM,PN$ 与椭圆 $E$ 的另一个交点分别为 $A,B$. 证明：直线 $AB$ 过定点.
\end{example}

\begin{example} % 入题思路（我搞复杂了），直线过顶点
（2022朝阳二模）已知椭圆 $C:\frac{x^2}{a^2}+\frac{y^2}{b^2}=1\ (a>b>0)$ 的一个顶点为 $P(0,1)$，离心率为 $\frac{\sqrt{2}}{2}$。\\
（1）求椭圆 $C$ 的方程；\\
（2）过点 $P$ 作斜率为 $k_1$ 的直线 $l_1$，交椭圆 $C$ 于另一点 $A$；过点 $P$ 作斜率为 $k_2\ (k_2\neq k_1)$ 的直线 $l_2$，交椭圆 $C$ 于另一点 $B$。若 $k_1k_2=1$，求证：直线 $AB$ 经过定点。
\end{example}
